% Generated draft: requires mathematical and source review.
\documentclass[11pt]{article}
\usepackage[margin=25mm]{geometry}
\usepackage{fontspec}
\setmainfont{Latin Modern Roman}
\usepackage{amsmath,amssymb,mathtools,unicode-math}
\setmathfont{Latin Modern Math}
\usepackage{enumitem,needspace,xurl,hyperref}
\setlength{\parindent}{0pt}
\setlength{\parskip}{5pt}
\setlength{\emergencystretch}{3em}
\newcommand{\N}{\mathbb{N}}
\newcommand{\Z}{\mathbb{Z}}
\newcommand{\Q}{\mathbb{Q}}
\newcommand{\R}{\mathbb{R}}
\newcommand{\C}{\mathbb{C}}
\newcommand{\F}{\mathbb{F}}
\newcommand{\A}{\mathbb{A}}
\newcommand{\PP}{\mathbb{P}}
\newcommand{\E}{\mathbb{E}}
\newcommand{\eps}{\varepsilon}
\newcommand{\dd}{\,\mathrm{d}}
\begin{document}
\section*{UnsolvedMath ID: 30001291}
\textbf{Source problem code:} catalogue entry
\subsection{Definitions and mathematical statement}
\textbf{Terminology and missing construction.} The supplied record names a ``simplex map'' and an ``invariant fundamental set'' but does not define either object. The following gives standard terminology only; it does not identify the map or impose hypotheses absent from the source.

A piecewise translation on a set \(X\subseteq\R^n\) is a map \(T:X\to X\) with a partition \(X=\bigsqcup_{i\in I}X_i\) and vectors \(v_i\in\R^n\) such that
\[
T(x)=x+v_i\qquad(x\in X_i).
\]
The supplied record does not specify the dimension, the set \(X\), the partition, the translation vectors, or the rule assigning boundary points to pieces.

For a full-rank lattice \(\Lambda\subset\R^n\), the quotient \(\R^n/\Lambda\) is a torus. A torus translation has the form
\[
\tau_a(x+\Lambda)=x+a+\Lambda
\]
for some \(a\in\R^n\). An exact fundamental set for \(\Lambda\) contains exactly one representative of each coset in \(\R^n/\Lambda\). Fundamental sets may instead be considered up to sets of measure zero or with boundary identifications; the supplied record does not specify the intended convention. It also does not specify whether invariance means \(T(F)=F\), forward invariance \(T(F)\subseteq F\), or another convention.

A subset is connected if it cannot be expressed as the union of two disjoint, nonempty subsets that are open in its relative topology. Under the usual convex-geometric convention, a polytope is the convex hull of finitely many points and is therefore convex. The source's phrase ``possibly non-convex polytope'' consequently requires a different, unspecified convention; a finite union of convex polytopes is one possible interpretation, but it is not selected here.

\textbf{Source question, pending specification.} Determine the relationship between the piecewise-translation and torus-translation roles of the intended simplex map. Determine whether its intended invariant fundamental set is connected and whether it has a possibly non-convex polytope shape in the source's intended sense.

The source supplies no construction, parameter range, lattice, translation, or quantifier specifying which simplex maps or fundamental sets are intended. It also does not specify whether the relationship sought is a factor relation, a conjugacy, or another correspondence. Thus the question above preserves the supplied target, but a uniquely quantified mathematical statement cannot be recovered from this record alone.
\subsection{Short English statement}
The question asks how an unspecified simplex map relates piecewise translations to translations of a torus, and whether its invariant fundamental set is connected and has a possibly non-convex polytope shape. The map, admissible parameters, fundamental-set convention, and intended meaning of shape must be recovered before a precise conjecture can be stated. The registry labels the problem open; this label is supplied curation, not an independently verified conclusion.
\subsection{Sources}
\begin{itemize}
\item Local problem registry, supplied record for permanent UnsolvedMath ID \(30001291\), source problem code \texttt{catalogue entry}, title ``Connections Between Simplex and Torus Translations,'' in the Oberwolfach Reports Open Problems collection. The record gives proposed year \(2009\), names no proposer, and labels the statement ``exact'' with the verification note that it is internally readable and retained byte-for-byte. These are supplied registry annotations.
\item Supplied report citation: \emph{Dynamische Systeme (2010)}, Oberwolfach Reports, DOI \texttt{10.4171/owr/2009/32}, pp.\ 1787--1842. \url{https://doi.org/10.4171/owr/2009/32}. The supplied background describes Oberwolfach Reports as documenting mathematical workshops and their open-problem discussions. No exact passage or map construction is supplied.
\item Supplied literature source labeled ``Adler--Nowicki--Swirszcz--Tresser'': \url{https://doi.org/10.1017/S0143385709000595}. The record provides no title or relevant theorem statement for this reference.
\item Literature assessment dated \texttt{2026-08-21}: the supplied curation labels the problem open and states that the report explicitly calls connectivity and shape of the invariant fundamental set open, and that later work does not provide a general resolution. It lists the two DOI references above as sources checked. This is supplied dated curation, not independently verified here; the record itself requests independent verification of time-sensitive status and cited claims.
\end{itemize}
\end{document}
